%%% FIELD thm-symbolic-null-classification-statement
The contraction--Möbius argument gives
\[
\mathcal N_3=\mathcal F_3,
\qquad
|\mathcal N_3|=1+9+12+6=28.
\]
More explicitly, fix a labeling \(U=\{a,b,c\}\) and abbreviate
\(\{a,b\}\) by \(ab\), and similarly for the other pairs.  In the
branch of a map \(g\in\mathcal N_3\) with \(g(U)=ab\), the four maps for which the remaining factor is
unresponsive alone, meaning \(g(\{c\})=\varnothing\), are indexed by
\((\rho_a,\rho_b)\in\{0,1\}^2\) and, for every \(z\subseteq U\), obey
\[
g(z)=
\begin{cases}
ab,&ab\subseteq z,\\
\{a\},&a\in z,\ \rho_a=1,\text{ and }ab\nsubseteq z,\\
\{b\},&b\in z,\ \rho_b=1,\text{ and }ab\nsubseteq z,\\
\varnothing,&\text{otherwise}.
\end{cases}
\]
The two maps for which the remaining factor is responsive alone,
meaning \(g(\{c\})=\{c\}\), have, in the coordinate order
\(\{a\},\{b\},\{c\},ab,ac,bc,U\), respectively the value vectors
\[
(\varnothing,\{b\},\{c\},ab,\{c\},\{b\},ab)
\quad\text{and}\quad
(\{a\},\varnothing,\{c\},ab,\{a\},\{c\},ab).
\]
Consequently the full-offer-pair branch contains exactly four maps
with \(c\) unresponsive alone and exactly two maps with \(c\)
responsive alone; together with the one full-offer-empty map and the
nine full-offer-singleton maps, these are exactly the four families in
\(\mathcal F_3\).

%%% FIELD thm-symbolic-null-classification-proof
Write \(\mathbf 1\{E\}\) for the indicator of an event or proposition
\(E\).  By one-sided receipt,
\(g(\varnothing)\subseteq\varnothing\), hence
\(g(\varnothing)=\varnothing\).  We first derive the consequence of
strict rational choice that will be used below.  If
\[
g(z)=r,
\qquad
r\subseteq w\subseteq z,
\]
then \(r\) is feasible from both offers \(z\) and \(w\).  By strict
rational choice, \(r\succ_g d\) for every other \(d\subseteq z\).
Every competitor feasible from \(w\) is also feasible from \(z\), so
\(r\) remains the unique \(\succ_g\)-maximal feasible bundle from
\(w\).  Therefore
\[
g(z)=r,\quad r\subseteq w\subseteq z
\quad\Longrightarrow\quad
g(w)=r. \tag{1}
\]

For distinct \(j,k\in U\), abbreviate \(\{j,k\}\) by \(jk\).
Expanding the definition of the Möbius signature on the three-element
Boolean lattice gives
\[
\sigma_U(g)=\mathbf 1\{g(U)=U\}, \tag{2}
\]
\[
\sigma_{jk}(g)
=\mathbf 1\{g(jk)=jk\}-\mathbf 1\{g(U)=jk\}, \tag{3}
\]
and
\[
\sigma_{\{j\}}(g)
=\mathbf 1\{g(\{j\})=\{j\}\}
-\sum_{k\in U\setminus\{j\}}
 \mathbf 1\{g(jk)=\{j\}\}
+\mathbf 1\{g(U)=\{j\}\}. \tag{4}
\]
Indeed, these are precisely the terms with supersets of cardinality
three, at least two, and at least one, respectively, and their signs
are \((-1)^{|z|-|t|}\), as prescribed by the definition.

We prove \(\mathcal N_3\subseteq\mathcal F_3\).  Let
\(g\in\mathcal N_3\).  Thus all its signatures vanish.  Equation
\((2)\) yields \(g(U)\ne U\).  One-sided receipt then leaves only the
three possibilities
\[
|g(U)|=0,
\qquad
|g(U)|=1,
\qquad
|g(U)|=2. \tag{5}
\]

If \(g(U)=\varnothing\), apply \((1)\) with
\(r=\varnothing\) and arbitrary \(w\subseteq U\).  It follows that
\(g(w)=\varnothing\) for every \(w\), which is family (i).

Suppose next that \(g(U)=\{a\}\), and write the other two factors as
\(b,c\).  Applying \((1)\) to every intermediate offer containing
\(a\) gives
\[
g(\{a\})=g(ab)=g(ac)=g(U)=\{a\}. \tag{6}
\]
Since \(g(U)\ne bc\), equations \((3)\) and
\(\sigma_{bc}(g)=0\) imply \(g(bc)\ne bc\).  One-sided receipt
therefore gives
\[
g(bc)\in\{\varnothing,\{b\},\{c\}\}. \tag{7}
\]
Substituting \((6)\) and \(g(U)=\{a\}\) into \((4)\), first for
\(j=b\) and then for \(j=c\), gives the two exact equations
\[
\mathbf 1\{g(\{b\})=\{b\}\}
=\mathbf 1\{g(bc)=\{b\}\},
\qquad
\mathbf 1\{g(\{c\})=\{c\}\}
=\mathbf 1\{g(bc)=\{c\}\}. \tag{8}
\]
Because a one-sided choice from a singleton is either that singleton
or the empty bundle, \((7)\)--\((8)\) determine the response map
completely: outside offers containing \(a\), there is either no
responsive singleton, responsive singleton \(b\), or responsive
singleton \(c\).  These are exactly the three family-(ii) maps for
the fixed primary factor \(a\).  There are three choices of \(a\), so
this branch has nine maps.

It remains to analyze \(g(U)=ab\), where \(c\) is the remaining
factor.  Contraction \((1)\) gives
\[
g(ab)=ab. \tag{9}
\]
Set \(x=g(ac)\) and \(y=g(bc)\).  Equations \((3)\), applied to
\(ac\) and \(bc\), show that \(x\ne ac\) and \(y\ne bc\), since the
corresponding full-offer indicators vanish.  Together with one-sided
receipt, this yields
\[
x\in\{\varnothing,\{a\},\{c\}\},
\qquad
y\in\{\varnothing,\{b\},\{c\}\}. \tag{10}
\]
Substitution of \((9)\)--\((10)\) and \(g(U)=ab\) into the three
instances of \((4)\) gives
\[
\mathbf 1\{g(\{a\})=\{a\}\}
=\mathbf 1\{x=\{a\}\}, \tag{11}
\]
\[
\mathbf 1\{g(\{b\})=\{b\}\}
=\mathbf 1\{y=\{b\}\}, \tag{12}
\]
and
\[
\mathbf 1\{g(\{c\})=\{c\}\}
=\mathbf 1\{x=\{c\}\}+\mathbf 1\{y=\{c\}\}. \tag{13}
\]

If \(g(\{c\})=\varnothing\), the left side of \((13)\) is zero;
both indicators on its right side must therefore be zero.  Hence
\[
x\in\{\varnothing,\{a\}\},
\qquad
y\in\{\varnothing,\{b\}\}. \tag{14}
\]
The two choices in \((14)\) are independent, while \((11)\) and
\((12)\) determine the singleton responses.  Thus there are exactly
\(2\cdot2=4\) such maps, precisely the family-(iii) maps indexed by
\((\rho_a,\rho_b)\in\{0,1\}^2\).

If instead \(g(\{c\})=\{c\}\), equation \((13)\) says that exactly
one of \(x,y\) equals \(\{c\}\).  Suppose first that
\(x=\{c\}\).  Since \(x\) is the strictly rational choice from
\(ac\), one has \(\{c\}\succ_g\varnothing\).  The bundle
\(\{c\}\) is feasible from \(bc\), so strict maximality makes
\(y=\varnothing\) impossible.  By \((10)\), and because
\(y\ne\{c\}\), necessarily
\[
y=\{b\}. \tag{15}
\]
Equations \((11)\)--\((12)\) then force
\(g(\{a\})=\varnothing\) and \(g(\{b\})=\{b\}\).  This is the
first displayed responsive-\(c\) value vector in the statement.  If
\(y=\{c\}\), the same argument with \(a,b\) exchanged forces
\[
x=\{a\},
\qquad
g(\{a\})=\{a\},
\qquad
g(\{b\})=\varnothing, \tag{16}
\]
which is the second vector.  Thus the responsive-\(c\) branch has
exactly two maps.  In particular, the putative seventh pattern is
excluded: when one of \(x,y\) equals \(\{c\}\), strict preference of
\(\{c\}\) to \(\varnothing\) prevents the other choice from being
empty, while \((13)\) prevents it from also being \(\{c\}\).
This proves that every null rational map belongs to one of the four
families.

For the reverse inclusion, we exhibit strict rationalizations and
verify every signature.  Since \(\mathcal B\) is finite, each of the
following displayed priority lists can be completed to a strict total
order by ordering all bundles in an indicated residual block
arbitrarily.  For family (i), rank \(\varnothing\) above every
nonempty bundle.  For family (ii), rank the primary singleton
\(\{a\}\) first, then the optional secondary singleton when it is
nonempty, then \(\varnothing\), and put every unlisted bundle below
\(\varnothing\).  For family (iii), rank \(ab\) first, then the
singletons among \(\{a\},\{b\}\) whose corresponding
\(\rho\)-value is one, then \(\varnothing\), and put every unlisted
bundle below \(\varnothing\).  Finally, the two family-(iv) maps are
rationalized respectively by
\[
ab\succ_g\{b\}\succ_g\{c\}\succ_g\varnothing\succ_g\{a\}
\quad\text{and}\quad
ab\succ_g\{a\}\succ_g\{c\}\succ_g\varnothing\succ_g\{b\},
\tag{17}
\]
with every unlisted bundle below \(\varnothing\).  Inspection of each
principal menu \(\{d:d\subseteq z\}\) gives exactly its defining
family rule.  Hence every map in \(\mathcal F_3\) is one-sided and
strictly rational.

It remains only to check nullity; this check also makes the use of all
eight signature equations explicit.  In family (i), no nonempty
bundle is ever selected, so all seven nonempty signatures are zero.
In family (ii), equation \((2)\) is zero, all pair signatures in
\((3)\) are zero, and
\[
\sigma_{\{a\}}(g)=1-1-1+1=0. \tag{18}
\]
If the optional secondary singleton is, say, \(\{b\}\), then
\[
\sigma_{\{b\}}(g)=1-0-1+0=0, \tag{19}
\]
while every term of \(\sigma_{\{c\}}(g)\) is zero; when there is no
secondary singleton all terms for both \(b,c\) are zero, and the
case with secondary singleton \(\{c\}\) is symmetric.  In family
(iii), equation \((3)\) gives
\(\sigma_{ab}(g)=1-1=0\) and gives zero for the other two pairs,
while
\[
\sigma_{\{a\}}(g)=\rho_a-0-\rho_a+0=0,
\qquad
\sigma_{\{b\}}(g)=\rho_b-0-\rho_b+0=0,
\qquad
\sigma_{\{c\}}(g)=0. \tag{20}
\]
Equation \((2)\) again gives zero.  For the first family-(iv) map in
the statement, the pair signatures are
\(\sigma_{ab}(g)=1-1=0\) and zero for \(ac,bc\), and the singleton
signatures are
\[
\sigma_{\{a\}}(g)=0,
\qquad
\sigma_{\{b\}}(g)=1-0-1+0=0,
\qquad
\sigma_{\{c\}}(g)=1-1-0+0=0. \tag{21}
\]
The second family-(iv) map is symmetric, and its triple signature is
zero by \((2)\).

For completeness, the empty-cell signature follows from the other
seven and one-sided receipt rather than requiring a further case
check.  Interchanging the two finite sums in the definition gives
\[
\begin{aligned}
\sum_{t\subseteq U}(-1)^{|t|}\sigma_t(g)
&=\sum_{z\subseteq U}(-1)^{|z|}
  \sum_{t\subseteq z}\mathbf 1\{g(z)=t\}\\
&=\sum_{z\subseteq U}(-1)^{|z|}\\
&=(1-1)^3=0.
\end{aligned} \tag{22}
\]
The inner indicator sum in the first line equals one because
one-sided receipt ensures that the unique value \(g(z)\) is a subset
of \(z\).  Since all nonempty signatures in
\((18)\)--\((21)\) and their stated symmetric cases vanish,
\((22)\) implies \(\sigma_\varnothing(g)=0\).  Thus
\(\mathcal F_3\subseteq\mathcal N_3\), completing the equality.

Finally, the empty, singleton, and pair cases are disjoint because
their full-offer choices have different cardinalities.  Within the
pair case, families (iii) and (iv) are separated by whether
\(g(\{c\})\) is empty or equals \(\{c\}\); within every branch the
displayed secondary or singleton choices determine the map.  Their counts are
\[
1,
\qquad
3\cdot3=9,
\qquad
3\cdot4=12,
\qquad
3\cdot2=6,
\]
respectively.  Therefore
\(|\mathcal N_3|=1+9+12+6=28\), as claimed.
